\documentclass[11pt,a4paper]{article}
\usepackage[margin=2.2cm]{geometry}
\usepackage{amsmath,amssymb,xcolor,enumitem,booktabs}
\usepackage[normalem]{ulem}
\usepackage[hidelinks]{hyperref}
\setlength{\parskip}{0.6em}\setlength{\parindent}{0pt}
\definecolor{newc}{HTML}{1f6f43}
\newcommand{\new}[1]{\textcolor{newc}{#1}}
\newcommand{\old}[1]{\textcolor{gray}{\sout{#1}}}
\newcommand{\cit}[1]{#1}
\newenvironment{greenblock}{\color{newc}}{}
\title{Proposed change, more of the model back in the Problem Formulation}
\date{17 November 2026, version 1}
\begin{document}\sloppy
\maketitle

\section*{How to read this}
Part~1 says why the chapter should grow. Part~2 goes through
\texttt{3.ProblemFormulation.tex} one subsection at a time and shows only the
paragraphs that change. Black text is there now and stays, as context.
\new{Green} is new text and \old{struck} text goes. Part~3 lists what changes
elsewhere, and Part~4 has the questions. Nothing is changed in the thesis yet.

\section*{Part 1. Why}

The chapter is about a third shorter than a month ago. Words without comments:

\begin{center}
\begin{tabular}{llr}
  \toprule
  date & what happened & words \\
  \midrule
  19 October & the longest it has been & 13{,}500 \\
  26 October & [A52] moves 25 definitions to Appendix C & 11{,}300 \\
  2 November & the readability rewrite trims the prose & 8{,}750 \\
  today & [A93], refills and figures added since & 9{,}330 \\
  \bottomrule
\end{tabular}
\end{center}

[A52] made the chapter lighter, from 51 displayed equations to 20. Your rule
of 23 October then said the formal definitions of the model never leave the
thesis, because the Problem Formulation is what is graded, and that the
statements in Appendix C go back when the thesis is reworked. This is that
rework.

Today the chapter uses symbols it never defines. The table of measures cites
the runtime $\kappa_T$. The observation function uses the section-based
policy. The gain test uses $F_{\mathrm{fix}}$ and the set of rules. The
research questions name the communication graph. A reader meets each of them
only in an appendix.

The proposal brings back eleven definitions, each in words first, with an
example on numbers the chapter already has. It adds a few reasons the
rewrite cut, and it fixes the direction to the target, which no longer
matches the code. It adds about 1{,}700 words, eleven displayed equations and
one table. Appendix C keeps the distance-label bound, the model against its
sources, the ties in the storage objective, the example agent's observation
and which crowding values a controller reads.

\section*{Part 2. The chapter as it would read}

% ------------------------------------------------------------------
\subsection*{2.1 The environment}

After ``\dots where an edge $(v,w)\in E$ is a move from $v$ to $w$.'':

\begin{greenblock}
The model uses a general graph and not a grid, because real warehouses have
one-way aisles, shortcuts and blocks of odd shape. A grid is one special
case. The positions are numbered once, $V=\{v_1,\dots,v_{|V|}\}$. The numbers
only break ties, for example between two shelves that are equally near.
\end{greenblock}

The last paragraph ends:

Every generated warehouse is checked so that an agent can reach every shelf,
workstation, dock and endpoint from wherever it stands. \new{This asks less
than that every position can reach every other, the condition of van
Knippenberg et al.\ (2021), and it is all the model needs.}

% ------------------------------------------------------------------
\subsection*{2.2 Agents and collisions}

In the paragraph on hard rules:

\dots so a deterministic conflict-resolution mechanism enforces the rules
after the agents have chosen. \new{A penalty in the reward would not be
enough, because a learned policy draws its actions at random and can still
draw a forbidden one.} The same mechanism is used for every controller, so
it is part of the environment rather than part of any policy.

% ------------------------------------------------------------------
\subsection*{2.3 Where the items are}

\textbf{How much arrives.} New, after the paragraph on units on their way:

\begin{greenblock}
How much stock arrives is set by the orders alone, so it is the same under
every storage rule. Each SKU counts the units ordered since its last refill,
from the moment a line is released. At a storage update it gets one refill
for every $Q$ units it has counted, and the rest carries over to the next
update. With $Q=10$, if 23 units of tea were ordered since the last refill,
two refills of tea arrive and 3 units carry over. This is the $(R,nQ)$
policy of inventory control (Axs\"ater 2015), with the reorder point one
refill below the starting stock.

Write $y^{\mathrm{ord}}_t(k)$ for the units of SKU $k$ ordered since its
last refill, with ``ord'' for ordered. Formally, the \emph{arrivals} are
\[
  y_t(k) = Q \left\lfloor \frac{y^{\mathrm{ord}}_t(k)}{Q} \right\rfloor ,
  \tag{arrivals}
\]
and the rest, $y^{\mathrm{ord}}_t(k)-y_t(k)$, fewer than $Q$ units, carries
over.

Every storage rule starts from the same stock $x_0\in\mathcal{X}$. Each SKU
starts alone on a shelf drawn at random for the seed, with one refill plus a
safety stock (Implementation chapter, instances).
\end{greenblock}

\textbf{What the rule knows.} After the list of the two numbers:

\begin{greenblock}
Both numbers look back over whole storage periods. The demand counts the
order lines for $k$ released in the last period and divides by its length
$\Delta$. If tea was ordered 30 times in a period of 100 steps,
$\hat\rho_t(\mathrm{tea})=0.3$. The standing still of a shelf is first
averaged over the picks from it that finished in one period. One period
holds only about eight picks per shelf in the 8-agent warehouse, too few for
a steady number, so $\hat\omega_t(v)$ averages the last five such periods. A
shelf with no finished pick yet reads 0.

Formally, write $P_s(v)$ for the picks from shelf $v$ that finished in
period $s$, and $S_t(v)$ for the last five periods up to $t$ in which any
did. The \emph{estimates} are
\[
  \begin{aligned}
  \hat\rho_t(k) &= \tfrac{1}{\Delta}\,\#\{\text{order lines for } k
    \text{ released in the last period}\}, \\
  \hat\omega_t(v) &= \frac{1}{|S_t(v)|}\sum_{s \in S_t(v)}
    \frac{1}{|P_s(v)|}\sum_{j \in P_s(v)} W_j ,
  \end{aligned}
  \tag{estimates}
\]
where $W_j$ is the blocked time of pick $\tau_j$, the third part of the
service-time split.
\end{greenblock}

\textbf{The rules compared.} New, after ``\dots and they do not search for the
smallest value (\dots).'' and before the example of arriving tea under fixed and demand storage:

\begin{greenblock}
A storage rule is any $F$ that meets the three conditions of the storage
update. Write $\mathcal{F}$ for all of them. The thesis compares three.
Fixed storage $F_{\mathrm{fix}}$ gives every SKU a home shelf at the start
and sends each refill back there, whatever the orders say. It is the
dedicated storage of a warehouse that never re-slots. Demand storage
$F_{\mathrm{dem}}$ and congestion-aware storage $F_{\mathrm{cng}}$ choose a
shelf at every update, as above. How each rule ranks the shelves is in
the storage rules of the Implementation chapter.
\end{greenblock}

% ------------------------------------------------------------------
\subsection*{2.4 Orders and tasks}

The pickup paragraph ends:

\dots Had the nearer shelf held only one unit, the line would take that one
and the other three from the second shelf, as two tasks. \old{This is the
\emph{pickup rule}~(eq:pickuprule).}

\begin{greenblock}
Formally, sort the shelves $v_1, v_2,\dots$ that hold free units of $k$ by
their distance $d_G(v_i,g)$ to the line's workstation $g$, and by vertex
number when equal. A line for $q$ units then takes
\[
  q^{(i)} = \min\Bigl(q - \sum_{l < i} q^{(l)},\; x_t(k, v_i) - r_t(k, v_i)\Bigr)
  \tag{pickup rule}
\]
units from shelf $v_i$, and every shelf with $q^{(i)}>0$ gives one task. In
the example with one unit on the nearer shelf, $q^{(1)}=\min(4,1)=1$ and
$q^{(2)}=\min(3,\,x_t-r_t)=3$, because the second shelf holds at least three.
\end{greenblock}

The congestion-aware paragraph ends:

\dots So congestion-aware storage acts slowly on where stock goes and quickly
on where it is taken from. \new{In the pickup rule its shelves are sorted by
$d_G(v_i,g)+\hat\omega_t(v_i)$ instead.}

% ------------------------------------------------------------------
\subsection*{2.5 A task's life}

Only a free agent can get a new task. An agent is free when it has no active
task. \new{Formally, the \emph{free-agent condition} is}
\begin{greenblock}
\[
  a_i \text{ is free at } t
  \iff
  \text{there is no } \tau_j \in \mathcal{B}_t \text{ with } \sigma(j)=i .
  \tag{free agent}
\]
\end{greenblock}
Agent $a_2$ is busy from the moment it gets $\tau_1$, also on its way to the
shelf, and is free again once $\tau_1$ is delivered.

% ------------------------------------------------------------------
\subsection*{2.6 Who decides}

\textbf{Zones.}

A section-based controller decides for the agents in one zone, and it only
sees that zone. The floor is cut into $p$ zones $Y_1,\dots,Y_p$ that do not
overlap and together cover every vertex an agent can stand on. \old{These form
the \emph{zone partition} $\mathcal{Y}$~(eq:partition).} \new{Formally, the
\emph{zone partition} is}
\begin{greenblock}
\[
  \mathcal{Y} = \{Y_1,\dots,Y_p\}, \qquad
  \bigcup_{q=1}^{p} Y_q = V, \qquad
  Y_q \cap Y_{q'} = \emptyset \ \text{ for } q \neq q' .
  \tag{zone partition}
\]
\end{greenblock}
When an agent crosses into a neighbouring zone, that zone takes it over.
Tasks are not handed over, because assignment is the same rule for the whole
warehouse.

\textbf{The section-based policy.}

A section-based controller does the same inside each zone. The agents in
zone $Y_q$ choose one at a time in the same priority order, each from what
the zone sees and the moves of the earlier agents in that zone. \old{This is
the \emph{section-based policy}~(eq:sectionpolicy), the centralised policy
with the whole warehouse replaced by the zone $s^{(q)}_t$, including the
strip just past its edge.} \new{Formally, write $s^{(q)}_t$ for what zone
$Y_q$ sees, its own part of the state and the strip just past its edge, and
$u^{(q)}_{\prec i}(t)$ for the moves of the agents before $a_i$ in that zone.
The \emph{section-based policy} is the centralised policy with $s_t$
replaced by $s^{(q)}_t$,}
\begin{greenblock}
\[
  u_i(t) \sim \pi^{\mathrm{sec}}_\theta\bigl(\cdot \mid s^{(q)}_t,\, u^{(q)}_{\prec i}(t),\, i\bigr),
  \qquad \ell_i(t) \in Y_q .
  \tag{section-based policy}
\]
\end{greenblock}
The zones choose at the same time and do not see each other's choices. \dots

% ------------------------------------------------------------------
\subsection*{2.7 What a decentralised agent sees}

\textbf{Direction to the target.} This part no longer matches the code. Two
things are missing.

First, the current target is defined only for an agent with a task. The code
also gives an idle agent a target, as [A72] decided, and \texttt{idle\_parking}
is on in every config. [A72] asked for two more cases in the current target
and a paragraph in this section. Neither is in the thesis now, and the
notation table still lists only the two task cases.

Second, the policy does not read the distance label itself. It reads the
difference from the agent's own vertex, capped at $\pm d_{\mathrm{obs}}$
([A78]), and the label divided by 50 ([A82]). That belongs in the
Implementation chapter, where it already is, but this section should say in
one sentence that the label reaches the policy in another form.

The proposed text:

Seeing nearby vertices does not tell the agent which way to go, because its
goal is usually outside its field of view. So every visible vertex gets a
label, its distance to the agent's current target. The target is the shelf
until the agent has picked up, and where the task delivers after that.
\new{An agent with no task keeps its own vertex as its target, so it stays
where it is. When it stands in another agent's way, on a shelf, workstation
or dock of an unfinished task, or on a vertex another agent was stopped from
moving onto, its target becomes the nearest free endpoint, until it arrives
there or gets a task. Without this, an idle agent would stay where its last
task ended, often on a dock or a workstation, and block the next agent that
needs it.}

Formally, \old{for agent $a_i$ serving task $\tau_j$, so $\sigma(j)=i$,} the
\emph{current target} \new{of agent $a_i$} is
\[
  q_i(t) =
  \begin{cases}
    s_j, & \text{\old{$t < t_j^{\mathrm{pickup}}$}\new{$a_i$ serves $\tau_j$ and $t < t_j^{\mathrm{pickup}}$}}, \\
    g_j, & \text{\old{$t \ge t_j^{\mathrm{pickup}}$}\new{$a_i$ serves $\tau_j$ and $t \ge t_j^{\mathrm{pickup}}$}}, \\
    \new{q^{\mathrm{park}}_i(t),} & \new{\text{$a_i$ is free and in the way}}, \\
    \new{\ell_i(t),} & \new{\text{$a_i$ is free otherwise}},
  \end{cases}
\]
\new{where $q^{\mathrm{park}}_i(t)$ is the nearest free endpoint, kept until
$a_i$ reaches it or gets a task.} The \emph{distance label} of a vertex $v$
is the travel distance from $v$ to that target, $\eta_i(v,t)=d_G(v,q_i(t))$.

The paragraph after the example ends:

The distance label is the shortest-path distance feature of van Knippenberg
et al.\ (2021). Because it is a distance and not just a ``closer or further''
flag, a policy can take a small step back for a bigger gain later. It
changes whenever the target changes. \new{The policy reads it in two forms,
as a difference from the label of the agent's own vertex and as the label
divided by a fixed number, so it stays small next to the other features
(Implementation chapter, the network).}

\textbf{How crowded the neighbourhood is.} The worked example stays and ends:

\dots so $\delta_1(t)=2/5=0.4$. Two fifths of the nearby space is taken.
\old{(eq:congestion) in Appendix~C states it formally.} \new{Formally, the
\emph{crowding feature} counts the other agents in the field of view and
divides by the other vertices there,}
\begin{greenblock}
\[
  \delta_i(t) =
  \frac{\bigl|\{a_j \in \mathcal{A}\setminus\{a_i\} :
                 \ell_j(t) \in V^{(d_{\mathrm{obs}})}_i(t)\}\bigr|}
       {\max\bigl\{1,\, |V^{(d_{\mathrm{obs}})}_i(t)| - 1\bigr\}} .
  \tag{crowding feature}
\]
The 1 keeps it defined when the field of view holds only the agent's own
vertex.
\end{greenblock}

(The sentence pointing to the appendix moves after the example, so the
number comes first and the formula confirms it.)

\textbf{Messages from nearby agents.}

\dots These links form the communication graph $G^{\mathcal{A}}_t$. Its
nodes are agents, and an edge $(a_i,a_j)$ means ``$a_i$ hears $a_j$''.
\new{Formally, $a_i$ hears $a_j$ when $a_j$ stands in its field of view,}
\begin{greenblock}
\[
  G^{\mathcal{A}}_t = (\mathcal{A}, E^{\mathcal{A}}_t), \qquad
  (a_i, a_j) \in E^{\mathcal{A}}_t \iff
  i \neq j \text{ and } d_G(\ell_i(t), \ell_j(t)) \le d_{\mathrm{obs}} .
  \tag{communication graph}
\]
\end{greenblock}
In the observation figure, \new{$(a_1,a_2)$ and $(a_1,a_3)$ are edges, so}
$a_2$ and $a_3$ can message $a_1$, but $a_4$, outside the field of view,
cannot.

% ------------------------------------------------------------------
\subsection*{2.8 The learning problem}

After the definition of the game and its elements:

\dots takes the state before the step, the joint action and the state after
it. \old{Appendix~C pairs each element with the object it restates.}
\new{Table~X pairs each element with where the chapter defined it.}

\begin{greenblock}
\begin{center}
\small
\begin{tabular}{llp{0.5\linewidth}}
  \toprule
  Element & Role & Defined as \\
  \midrule
  $\mathcal{A}$   & agents            & the fleet $\{a_1,\dots,a_m\}$, Agents and collisions \\
  $\mathcal{S}$   & states            & where agents stand, the task stages, the stock $x_t$, and the priority order, drawn once per episode \\
  $\mathcal{U}_i$ & actions           & wait plus $d_{\max}$ numbered moves, see below \\
  $P$             & transition rule   & the move rule with conflict resolution, the shared assignment rule and the storage update. Random only through the orders \\
  $\mathcal{O}_i$, $O_i$ & observations & the observation function below, one case per controller \\
  $R_i$           & reward            & left open here, given in the Implementation chapter \\
  $\gamma$        & discount          & a training value, Implementation chapter \\
  \bottomrule
\end{tabular}
\end{center}

A stochastic game gives every agent one fixed list of actions, but the legal
moves depend on where the agent stands. So the list is as long as the
busiest vertex needs, one wait plus one move for each outgoing edge of the
vertex with the most. Move number $k$ takes the $k$-th outgoing edge of
wherever the agent stands, in vertex-number order. At a vertex with fewer
edges the extra moves are switched off before the policy chooses
(Implementation chapter, the network), so what remains is exactly $\mathcal{U}(v)$. In a warehouse
whose busiest junction has four exits, every agent has five actions, and an
agent in a one-way aisle has two of them switched on. Formally,
\[
  \mathcal{U}_i = \{\mathrm{wait},\, \mathrm{move}_1, \dots, \mathrm{move}_{d_{\max}}\},
  \qquad d_{\max} = \max_{v \in V} |N^{+}(v)| .
  \tag{action list}
\]
\end{greenblock}

% ------------------------------------------------------------------
\subsection*{2.9 How a run is scored}

\textbf{Reported measures.}

Two measures are reported next to the result. The mean service time of
put-aways shows what demand and congestion-aware storage cost in agent time,
since a rule that puts stock near the workstations can lengthen the
put-aways. \new{It is the mean of $\zeta_j$ over the put-aways finished by
$T$. Fixed storage pays it too, because its refills also go to a shelf.} The
mean runtime $\kappa_T$ shows how long the controller takes to decide.
\old{Both are defined in Appendix~C.} \new{Write $\kappa(t)$ for the
wall-clock time the controller needs to produce the joint action from $s_t$,
assignment and routing together. Formally, the \emph{mean runtime} is}
\begin{greenblock}
\[
  \kappa_T = \frac{1}{T}\sum_{t=0}^{T-1} \kappa(t) ,
  \tag{runtime}
\]
measured in seconds, as Ma et al.\ (2017) report it.
\end{greenblock}

\textbf{The decision problem.} New, before ``The research question,
formally'':

\begin{greenblock}
Put together, the problem asks which pair of storage rule and controller
gives the shortest order service time, among the pairs whose fleet keeps up.
The run must keep up. No shelf may hold more than its room. Every move must
be legal and free of collisions. A pick may only take stock that is on its
shelf and not promised. Formally, the best pair is
\[
  (F^{\star}, \pi^{\star}) \in
      \arg\min_{F \in \mathcal{F},\, \pi \in \Pi} \mathbb{E}\bigl[\bar\zeta_T(F,\pi)\bigr],
  \tag{decision problem}
\]
with the expectation over the order stream and the seeds, over the pairs
that meet the four conditions. In the same order, they are
\begin{align*}
  &\text{the keep-up check} && \text{the run keeps up,}\\
  &\text{the room condition} && \text{no shelf over its room,}\\
  &\text{the move, vertex and swap rules} && \text{legal moves without collisions,}\\
  &\text{the stock condition} && \text{stock on the shelf and free.}
\end{align*}

The conditions do different jobs. The move, vertex and swap rules hold for
every controller, because conflict resolution enforces them, so they remove
no controller. They are listed so that a solution describes a run that can
happen. The room and stock conditions limit which rules and which tasks are
allowed. The keep-up check limits which pairs count at all.

The thesis does not search for the best pair. It compares three rules under
three controllers, and the gain test below asks whether a rule beats fixed
storage under one controller.
\end{greenblock}

% ------------------------------------------------------------------
\subsection*{2.10 Scope and assumptions}

New item in the list (see Question~4):

\begin{greenblock}
\begin{description}
  \item[Docks without limit] A dock takes any number of finished orders and
    refills. Only the vertex rule limits it, one agent at a time. Ma et al.\
    (2017) put no limit on pickup or delivery locations either.
\end{description}
\end{greenblock}

The two paragraphs after the list:

Tasks carry no deadline. A late delivery only lengthens service time, and no
single task can be said to have failed. This is narrower than
time-constrained formulations such as van Knippenberg et al.\ (2021), where a
missed deadline is infeasible or penalised. \new{Adding deadlines would put a
deadline in every task, and either forbid a late task or count late tasks
next to the result.}

Assignment is also permanent once made, so an agent serving a task is not
re-tasked before it finishes. Ma et al.\ (2017) allow re-tasking until the
agent reaches the pickup vertex. \old{Appendix~C says what adding deadlines or
re-tasking would change in the model.} \new{Adding it would let the
assignment rule change the agent of an active task whose agent has not yet
reached the shelf. Either change alters what counts as a solution, so it
would change this chapter, not the method.}

\section*{Part 3. What changes elsewhere}

\begin{itemize}[leftmargin=1.4em]
  \item \textbf{Appendix C.} ``The arrivals'', ``The starting stock'', the
    game table with the action list, the runtime, the decision problem,
    ``Deadlines'' and ``Re-tasking'' leave it. ``Definitions in full'' keeps
    only the Implementation definitions (trip distance, halo, features, masks,
    one-step cost). ``Capacity at the docks'' and its TODO go. The labels
    keep their names, so every reference still works, and
    \texttt{tools/sync\_xrefs.py --build} is run afterwards.
  \item \textbf{Notation.} All symbols are already in the table except
    $q^{\mathrm{park}}_i(t)$, which is added. $q_i(t)$ gets the two idle
    cases. $y$ is not used for the parking target, because $y_t(k)$ is the
    arrivals.
  \item \textbf{Code.} No change. Each new text was checked against
    \texttt{origin/main}. The refills (\texttt{storage/arrival.py}), the last
    five readings (\texttt{storage\_estimates.py}), the crowding denominator
    (\texttt{observation.py}), the idle target (\texttt{reward\_function.py},
    \texttt{current\_target}) and \texttt{idle\_parking: true} in every config.
  \item \textbf{GAPS.} One entry, [A97], for the move back and for the idle
    target. The idle target is a catch-up of [A72], not a new decision.
\end{itemize}

\section*{Part 4. Questions}

\begin{enumerate}[leftmargin=1.4em]
  \item \textbf{How many equations.} [A52] cut the chapter to 20 displayed
    equations, and it now has 34. This adds eleven. Keep all of them
    displayed, or write the short ones (free agent, zone partition, runtime)
    inline in a sentence?
  \item \textbf{The estimates.} The display folds the mean per period into
    one line with $W_j$. Keep that, or show the per-period mean $b_s(v)$ as
    its own line, as Appendix C does now?
  \item \textbf{The idle target.} Put the two idle cases in the current
    target, as proposed, or give them their own short paragraph after it and
    keep the equation to the two task cases?
  \item \textbf{Docks without limit.} The capacity of the docks is an open
    TODO in Appendix C. Is ``a dock takes any number, only the vertex rule
    limits it'' right?
  \item \textbf{Well-formed warehouses.} The three sentences on Token
    Passing's well-formed condition are commented out in ``Agents and
    collisions''. Bring them back in ``The environment'', or leave them in the
    appendix, since they name a method before the Implementation chapter?
  \item \textbf{What stays in Appendix C.} The distance-label bound, the
    model against its sources, the ties, the example agent's observation and
    which crowding values a controller reads. Agree?
\end{enumerate}

\end{document}
